% Korean edition of figs/schedule.tex.
\fig[wide]{schedule}{래퍼는 08:50 전에 준비를 마치고, 장중에는 폴링과 매매만 한다}[개장 때 한 번 리밸런싱하는 전략의 거래일]{%
{\tabfont\begin{tabularx}{\linewidth}{>{\RaggedRight\arraybackslash}p{20mm} >{\RaggedRight\arraybackslash}p{26mm} >{\RaggedRight\arraybackslash}X}
\thd{시각 (KST)} & \thd{단계} & \thd{래퍼 동작}\\
07:10 & 시작 & 저장된 쿠키 묶음으로 토큰을 갱신하고, 실패하면 로그인. 포트폴리오 선택. \code{currBizDay}가 KST 날짜와 같아질 때까지 주문 거부\\\rowsep
07:15 & 불러오기 & 실시간 피드를 열고 스냅샷을 받음. 종목별 섹터를 위해 \code{Universe}를 불러옴. 서버에서 그날 주문 건수를 셈\\\rowsep
08:00\textasciitilde 08:50 & 계획(호출자) & 호출자가 자체 가격으로 목표 비중을 계산. \code{plan()}이 유니버스, 합계, 종목·섹터 한도, 1주 비중과 대조해 주문으로 바꿈. \code{guidelines()}가 회전율과 소형주 합계를 보고\\\rowsep
08:50 & 제출 & \code{rebalance()}가 서버에서 그날 주문을 다시 세고 하루 한도 80건을 확인한 뒤, 매도부터 최소 1\,s 간격으로 보냄. 각 주문은 피드의 \code{ords}에 나타나면 확인됨\\\rowsep
09:00\textasciitilde 15:20 & 장중 & 3\,s마다 피드를 폴링. 720\,s 뒤 첫 REST 호출이 토큰을 갱신하고, 피드는 재연결할 때만 토큰을 요청. 호출자의 요청에 따라 주문 정지\\
\end{tabularx}}}{시각은 호출자가 정한다. 3\,s 폴링을 빼면 래퍼는 호출될 때만 동작한다. 08:30에 넣은 주문은 08:33에 \code{Working}이 됐고, 장 시작인 09:00에 체결되기 시작해 09:08:07에 \code{Done}이 됐다(\ref{sec:orders}절).}{DART Quant Team}
